\section*{GitHub-Tutorial}
\label{ch:GitHub}

In diesem Tutorial lernen Sie Schritt für Schritt, wie Sie:
\begin{itemize}
  \item einen \textbf{GitHub-Account} erstellen,
  \item ein neues \textbf{Repository} anlegen,
  \item Ihren Rechner via \textbf{SSH} mit GitHub verbinden,
  \item und ein bestehendes Python-Projekt mit \textbf{Git} (\lstinline|add|, \lstinline|commit|, \lstinline|push|) hochladen.
\end{itemize}

Git bietet folgende Haupt-Vorteile gegenüber dem Arbeiten auf dem lokalen Computer:
\begin{itemize}
  \item \textbf{Versionskontrolle}: Änderungen werden protokolliert, und frühere Versionen können leicht wiederhergestellt werden.
  \item \textbf{Zusammenarbeit}: Mehrere Personen können gleichzeitig an einem Projekt arbeiten, ohne sich gegenseitig zu stören.
  \item \textbf{Backup}: Der Code ist sicher in der Cloud gespeichert und kann von überall abgerufen werden.
  \item \textbf{Automatisierte Pipelines}: Möglichkeit, bei jeder neuen Version einen automatischen Test- und Deployment-Prozess zu starten, beispielsweise, um eine Dash-App auf einem Server zu aktualisieren.
\end{itemize}

\subsection*{Git und VS Code vorbereiten}
Installieren Sie Git mit folgendem Befehl:  

\begin{lstlisting}[style=bash]
# macOS
brew install git

# Linux (Ubuntu/Debian)
sudo apt install git
\end{lstlisting}

Für Windows: Laden Sie Git von \url{https://git-scm.com/download/win} herunter und installieren Sie es.

\subsection*{GitHub-Account erstellen}
Öffnen Sie \url{https://github.com} und klicken Sie auf \textbf{Sign up}. Folgen Sie den Schritten, um einen Account anzulegen.  

\subsection*{SSH-Schlüssel erstellen und hinterlegen}
Erzeugen Sie ein Schlüsselpaar, laden Sie es in den SSH-Agent und fügen Sie den Public Key bei GitHub ein.

Führen Sie folgenden Code auf Ihrem Computer aus (Terminal auf macOS oder Linux, bzw. Git Bash auf Windows):
\begin{lstlisting}[style=bash]
ssh-keygen -t ed25519 -C "ihre_email@schule.ch"
eval "$(ssh-agent -s)"
ssh-add ~/.ssh/id_ed25519
pbcopy < ~/.ssh/id_ed25519.pub  # Public Key in Zwischenablage kopieren
\end{lstlisting}


Auf GitHub: \textbf{Settings} $\rightarrow$ \textbf{SSH and GPG keys} $\rightarrow$ \textbf{New SSH key}.  

Testen Sie die Verbindung:  
\begin{lstlisting}[style=bash]
ssh -T git@github.com
\end{lstlisting}

\subsection*{Git konfigurieren}
Damit bei jedem commit auf Git die richtigen Angaben verwendet werden, konfigurieren Sie Git mit Ihrem Namen und Ihrer E-Mail-Adresse. Zusätzlich stellen Sie den Standard-Branch auf \lstinline|main| ein:

\begin{lstlisting}[style=bash]
git config --global user.name "Vorname Nachname"
git config --global user.email "vorname.nachname@stud.edu.zh.ch"
git config --global init.defaultBranch main
\end{lstlisting}

\subsection*{Neues Repository auf GitHub}
Erstellen Sie über das \enquote{+} oben rechts ein \textbf{New repository}, geben Sie einen Namen (z.\,B. \lstinline|mein-python-projekt|) und wählen Sie \enquote{Public}.



\subsection*{Projekt vorbereiten}
Wechseln Sie in den Ordner Ihres Projekts:  
\begin{lstlisting}[style=bash]
cd /pfad/zu/projekt
\end{lstlisting}

Erstellen Sie eine \lstinline|.gitignore|-Datei für Python:
\begin{lstlisting}
__pycache__/
*.pyc
.venv/
.DS_Store
.vscode/
\end{lstlisting}

\subsection*{Initialisieren, Committen, Pushen}
\begin{lstlisting}[style=bash]
git init
git add .
git commit -m "Erster commit!"

git remote add origin git@github.com:<user>/<repo>.git
git push -u origin main
\end{lstlisting}

\subsection*{Arbeiten mit VS Code}
Öffnen Sie den Projektordner in VS Code.  
Verwenden Sie die \textbf{Source Control}-Ansicht (Git-Symbol links), um Änderungen zu \textbf{stagen}, zu \textbf{committen} und zu \textbf{pushen}.

\subsection*{Nützliche Befehle}
\begin{lstlisting}[style=bash]
# Änderungen abrufen
git pull

# Remote-URL prüfen
git remote -v
\end{lstlisting}
